\scanpage{24}
\deleted{Si $x\in X\setminus Y$, $\exists\alpha$ tel que $f_\alpha(x)\ne0$, et}

définissent $Y$ comme l'ens des zéros communs. \deleted{de $X$} : ces fonctions, de plus elles séparent les pts de $X\setminus Y$, car si $x,y$ sont des pts distincts, \deleted{$\exists\beta$ tel que $f_\alpha(x)\ne0$ $g_\beta(x)\ne g_\beta(y)$ d'où}

Si les $f_\alpha$ \unsure{ne} suffisent pas à les séparer, montrons que les $f_\alpha g_\beta$ les séparent : en effet, $\exists\beta$ tel que $g_\beta(x)\ne g_\beta(y)$ \deleted{d'où} et $\alpha$ tel que $\lambda=f_\alpha(x)\ne0$, donc $f_\alpha(y)$ $(=f_\alpha(x)=\lambda\text{ par hyp.})\ne0$, d'où
\[
f_\alpha\underbrace{g_\beta(x)}_{\lambda}=\lambda g_\beta(x)\ne f_\alpha g_\beta(y)=\lambda g_\beta(y)
\]

Donc :

\uline{Proposition}. Si on a M1 : M\unsure{4|7} et M11, alors pour tt partie \deleted{\unsure{telle|fermée}} \inserted{modérée} $Y$ d'un esp. modéré $X$, qui est ens des zéros communs d'une famille finie de fonctions modérées sur $X$, on peut trouver une telle famille qui sépare les pts de $X$ (donc $X/Y$ est un espace modéré).

\mspage{23}
